\section{Pipelining}
Pipelining is an implementation of instruction execution, splitting each
instruction into execution stages, enabling the concurrent execution of
multiple instructions (at different execution stages). Pipelining
\textbf{increases the throughput of the entire workload} without decreasing
latency of any single task, but is \textit{limited by the slowest pipeline
stage} and requires \textit{stalling for dependencies}.\\

In MIPS, each instruction is split into \textbf{five execution stages}:
\begin{itemize}
  \item IF: Instruction Fetch
  \item ID: Instruction Decode and Register Read
  \item EX: Execute an Operation or Address Calculation
  \item MEM: Operand Access in Data Memory
  \item WB: Write Back the Result into a Register
\end{itemize}

\subsection{Implementation: Datapath}
In each cycle, we have the capabilities to execute an instruction in each
pipeline stage. That is, we may execute an instruction currently in the WB
stage, an instruction in the MEM stage, and so on. However, this necessitates
proper transferral and storage of data that is required by the instruction
across the pipeline stages.

\paragraph{Pipeline Registers}
To store the data use by the \textit{same instruction} in later pipeline
stages, the datapath stores additional registers for each of the following
stages: IF/ID, ID/EX, EX/MEM, MEM/WB.
\begin{remark}
  Instructions that do not require certain execution cycles (eg. MEM for ALU
  instructions) still have to go through the pipeline stage.
\end{remark}

\begin{figure}[h!]
  \centering
  \includegraphics[width=0.8\textwidth]{pipeline_data}
\end{figure}

\paragraph{IF Stage} At the end of a cycle, the IF/ID registers stores:
\begin{itemize}
  \item Instruction read from $InstructionMemory[PC]$
  \item $PC+4$
\end{itemize}

\begin{multicols}{2}
\paragraph{ID Stage} IF/ID registers supplies:
\begin{itemize}
  \item Read Registers for Register File
  \item $16$-bit offset to be sign-extended
  \item $PC+4$
\end{itemize}
ID/EX registers stores:
\begin{itemize}
  \item Read Data from Register File
  \item $32$-bit immediate value
  \item $PC+4$
\end{itemize}
\begin{itemize}
  \item[\vspace{\fill}]
\end{itemize}
\begin{itemize}
  \item Write Register
\end{itemize}
\end{multicols}

\begin{multicols}{2}
\paragraph{EX Stage} ID/EX registers supplies:
\begin{itemize}
  \item Read Data from Register File
  \item $32$-bit immediate value
  \item $PC+4$
  \item Write Register
  \item[\vspace{\fill}]
\end{itemize}
\columnbreak%
EX/MEM registers stores:
\begin{itemize}
  \item $\left( PC+4 \right) + \left( Immediate \times 4 \right) $
  \item ALU Output
  \item \verb|isZero?| Signal
  \item Read Data 2 from Register File
  \item Write Register
\end{itemize}
\end{multicols}

\begin{multicols}{2}
\paragraph{MEM Stage} EX/MEM registers supplies:
\begin{itemize}
  \item $\left( PC+4 \right) + \left( Immediate \times 4 \right) $
  \item ALU Output
  \item \verb|isZero?| Signal
  \item Read Data 2 from Register File
  \item Write Register
\end{itemize}
\columnbreak
MEM/WB registers stores:
\begin{itemize}
  \item ALU Output
  \item Read Data from Data Memory
  \item Write Register
  \item[\vspace{\fill}]
  \item[\vspace{\fill}]
\end{itemize}
\end{multicols}

\paragraph{WB Stage} MEM/WB registers supplies:
\begin{itemize}
  \item ALU Output
  \item Read Data from Data Memory
  \item Write Register
\end{itemize}
At the end of the cycle, result (if any) is written back to register file with
the supplied write register.

\subsection{Implementation: Control}
\begin{figure}[h!]
  \centering
  \includegraphics[width=0.8\textwidth]{pipeline_control2}
\end{figure}
The pipeline implementation of control signals is similar to that of a
single-cycle datapath. Each control signal is now associated with a particular
pipeline stage. Since the control signals are generated in the \textbf{ID
stage}, they are propagated through the pipeline registers until utilised.
\begin{figure}[h!]
  \centering
  \includegraphics[width=0.7\textwidth]{pipeline_control}
\end{figure}

\subsection{Comparison of Performances}
\begin{figure}[h!]
  \centering
  \includegraphics[width=0.5\textwidth]{comparison}
\end{figure}
\paragraph{Single-Cycle Processor}
With a cycle time of
\[CT_{seq}=\max{\left( \sum_{k=1}^{N} T_k \right) }\]
where $T_k$ is the time of operation in stage $k$ and $N$ is the number of
stages, the total execution time for $I$ instructions is
\[Time_{seq}=I\times CT_{seq}.\] 

\paragraph{Multi-Cycle Processor}
With a cycle time of
\[CT_{multi}=\max{\left( T_k \right) }\]
where $T_k$ is the time of operation in stage $k$, the total execution time
for  $I$ instructions is
\[Time_{multi}=I\times Average\ Cycles\ Per\ Instruction\times CT_{multi}.\]

\paragraph{Pipelining Processor}
With a cycle time of
\[CT_{pipeline}=\max{\left( T_k \right) }+T_d\]
where $T_k$ is the time of operation in stage $k$, and $T_d$ is the overhead
for pipelining, the total number of cycles needed for $I$ instructions is
$I+N-1$ (where $N$ is the number of stages), and the total execution time is
\[Time_{pipeline}=\left( I+N-1 \right)\times CT_{pipeline}.\] 
In an ideal case of pipelining, where:
\begin{itemize}
  \item Every stage takes the same amount of time. $\sum_{k=1}^{N} T_k=N\times T_k$
  \item No pipeline overhead. $T_d\to 0$
  \item Number of instructions is much larger than number of stages. $I >> N$
\end{itemize}
The pipeline processor can gain $N$ times speedup compared to a single-cycle
processor.

\subsection{Pipeline Hazards}
Pipeline hazards are issues that prevent the next instruction from immediately
being ``pumped'' into the pipeline following the previous instruction. These
can arise from various problems, such as \textbf{structural hazards}
(simultaneous use of a hardware resource), or instruction dependencies such as
\textbf{data hazards} (data dependencies between instructions) and
\textbf{control hazards} (change in program flow).

\subsubsection{Structural Hazards}
To solve structural hazards, we ensure that processes do not access or use the
same ``part'' of the hardware resource at the same time. This can be done in
several ways, such as:
\begin{enumerate}
  \item \textbf{Stalling} the pipeline for a cycle effectively offsets the next
    instructions by 1.
  \item Dividing the hardware resource (eg. Separating instruction memory and
    data memory)
\end{enumerate}
\begin{example}[Register Access]
  In both the EX and WB stages of the pipeline, instructions access the
  register file to read/write into the registers. To prevent problems in the
  case of accessing the same register, we \textit{split the cycle} into a write
  phase, and a read phase. Instructions are allowed to write into the register
  only during the first half of the cycle, and read from the register only
  during the second half of the cycle. This can be be done because registers
  comprises fast memory devices.
\end{example}

\subsubsection{Data Hazards}
Instructions can have various relationships that prevent pipeline execution.
When different instructions access the same register, we have \textbf{data
dependencies} between the instructions.
\begin{definition}[Read-After-Write]
  RAW occurs when a later instruction reads from the destination register which
  is written by an earlier instruction.
\end{definition}
\begin{figure}[h!]
  \centering
  \includegraphics[width=0.6\textwidth]{data_hazard}
\end{figure}
\begin{example}[RAW Data Hazard]
  In the code fragment, the multiple reads of $\$2$ after its initial write
  causes data hazards in the pipeline.
  \begin{lstlisting}
    sub $2, $1, $3
    and $12, $2, $5
    or $13, $6, $2
    add $14, $2, $2
    sw $15, 100($2)
  \end{lstlisting}
\end{example}
\paragraph{Forwarding} If the result is produced before it is used, that is,
the result exists in the pipeline registers, we may forward the result to any
trailing instructions from the pipeline registers, bypassing the read
registers.
\begin{figure}[h!]
  \centering
  \includegraphics[width=0.6\textwidth]{forwarding}
\end{figure}

\paragraph{Stalling} If the result is produced only in the MEM stage (in a LOAD
instruction), the result is required before it is produced, and we have to stall
the pipeline 
\begin{figure}[h!]
  \centering
  \includegraphics[width=0.5\textwidth]{stalling}
\end{figure}

\begin{remark}
  RAW is also known as true data dependency, as cases of Write-After-Write,
  Write-After-Read, Read-After-Read do not impose a data hazard.
\end{remark}

\subsubsection{Control Hazards}
\begin{definition}[Control Dependency]
  An instruction $j$ is control dependent on $i$ if $i$ controls whether $j$
  executes.
\end{definition}
In particular, the branch and jump instructions introduce control dependencies,
where the branch decision is made only in the MEM stage. The pipeline wrongly
executes 3 instructions and has to be flushed. The 3-cycle penalty can be
mitigated with \textbf{early branch resolution}, \textbf{branch prediction},
and \textbf{delayed branching}.

\paragraph{Early Branch Resolution}
By moving the branch target address calculation and register comparison, we can
make the decision in the ID stage instead of MEM (by creating a semi-ALU to
compare register values). This effectively reduces the delay from 3 to 1 clock
cycle.
\begin{figure}[h!]
  \centering
  \includegraphics[width=0.8\textwidth]{earlyBranch}
\end{figure}

However, if the registers involved in the comparison is data dependent on
preceding instructions, forwarding is required, introducing further stalls.

\begin{figure}[h!]
  \centering
  \includegraphics[width=0.49\textwidth]{earlyBranch2}
  \includegraphics[width=0.5\textwidth]{earlyBranch3}
\end{figure}

\paragraph{Branch Prediction}
We adopt the simple branch prediction scheme of assuming the branch not taken,
and continue ``pumping'' the successor instructions into the pipeline. If the
assumption is true, zero stall is required. If the assumption is false (the
branch is taken), then the piped instructions should not be executed, and are
flushed from the pipeline.

\begin{figure}[h!]
  \centering
  \includegraphics[width=0.6\textwidth]{branchPred}
\end{figure}

If early branching is adopted, this only causes a 1 cycle delay. Otherwise, 3
instructions would have been ``pumped'' and introduces a 3 cycle delay.

\paragraph{Delayed Branch}
To make use of the cycles that a branch might stall on, we may move non-control
dependent instructions that are executed regardless of the branch outcome.

\begin{multicols}{2}
  \begin{lstlisting}
    or  $8, $9, $10 // Non-control dependent
    add $1, $2, $3
    sub $4, $5, $6
    beq $1, $4, Exit
    xor $10, $1, $11

  Exit:
  \end{lstlisting}
  \begin{lstlisting}
    add $1, $2, $3
    sub $4, $5, $6
    beq $1, $4, Exit
    or  $8, $9, $10 // Branch-Delay Slot
    xor $10, $1, $11

  Exit:
  \end{lstlisting}
\end{multicols}

In MIPS, we have 1 branch-delay slot (with early branching). If there is no
such instruction that can be moved while preserving program correctness, a
\verb|nop| instruction is added.
